\documentclass[12pt,a4paper]{letter}
\usepackage[utf8]{inputenc}
\usepackage[T1]{fontenc}
\usepackage{lmodern}
\usepackage[margin=1in]{geometry}

\signature{Ian Todd\\Sydney Medical School\\University of Sydney}
\address{Ian Todd\\Sydney Medical School\\University of Sydney\\Sydney, NSW, Australia\\itod2305@uni.sydney.edu.au}

\begin{document}

\begin{letter}{Editor\\BioSystems}

\opening{Dear Professor Igamberdiev,}

I am writing to submit the manuscript ``Information Geometry of the Quantum-Classical Boundary in Biological Information Processing'' for consideration in \textit{BioSystems}.

This paper addresses a question central to BioSystems' scope: how does the quantum-to-classical transition constrain biological information processing? We prove that decoherence is a dimensional collapse in the quantum Fisher information manifold---the dephasing channel maps a $(d^2{-}1)$-dimensional quantum state manifold onto the $(d{-}1)$-dimensional classical probability simplex---and characterize this collapse geometrically using Bures principal angles. The mathematical results (Theorems~3.1--3.7) provide exact, closed-form characterizations: the Bures separation angle at every full-rank qubit state, and a global orthogonality theorem establishing that for $d \geq 3$, diagonal and off-diagonal tangent subspaces are Bures-orthogonal if and only if the state is diagonal.

The biological applications connect directly to your program on quantum coherence and thermodynamic limits in living systems. We apply the framework to three systems spanning 12 orders of magnitude in dephasing:

\begin{enumerate}
\item \textbf{Photosynthetic complexes} (FMO, $d = 7$; PE545, $d = 8$): At their ENAQT-optimal dephasing rates, all principal angles exceed $87^\circ$ (FMO) and $89^\circ$ (PE545)---biologically optimal transport occurs in a regime that is geometrically close to classical.

\item \textbf{Neural oscillations} ($d = 10$): Cortical gamma-band assemblies are ${\sim}12$ orders of magnitude past the boundary, geometrically indistinguishable from the classical simplex.
\end{enumerate}

An appendix extends the analysis to an illustrative molecular-oscillator model of chromatin dynamics ($d = 8$), exploring whether molecular biology generically sits near the quantum-classical boundary because $kT \sim J_{\max}$ at molecular energy scales. The energy-scale argument is suggestive, but a sensitivity sweep reveals that the geometric position is sensitive to coupling strength, motivating future work with experimentally determined Hamiltonians.

We connect these empirical results to the proved thermodynamic fact that coherences store free energy $k_BT \cdot C_r(\rho)$, suggesting a cost-benefit interpretation: the ENAQT regime represents a sweet spot where extra dimensional capacity from partial coherence comes at modest thermodynamic cost. This interpretation complements your work on temporal organization and thermodynamic limits in biological systems.

All claims are explicitly tagged as \emph{proved} (theorems), \emph{empirical} (computed from validated or physically motivated models), or \emph{hypothesis} (interpretation connecting the two). The full simulation code is available in the repository.

This paper builds on two previous BioSystems publications: ``Limits of Falsifiability in Complex Biological Systems'' (BioSystems 258, 105608, 2025) and ``Coherence Time in Biological Oscillator Assemblies Bounds the Rate of State Registration'' (accepted, 2026). The present work provides the quantum-mechanical foundation that those papers assumed: a geometric characterization of where the quantum-classical boundary lies and what biology can exploit on either side.

The manuscript has not been submitted elsewhere and is not under consideration at any other journal. There are no conflicts of interest.

\closing{Sincerely,}

\end{letter}
\end{document}
